% SOURCE_AUTHORITY: sga5_fr_workpass.tex, lines 8557--8586 (LF-normalized unit).
% SOURCE_SHA256: 791F4EFFC5E02832D5D77ED03518C8156D6F07E4C8238B03545DB93D883FBB28
\subsubsection*{1.4. Diversas categorías construidas a partir de la categoría de los \(\mathbb Z_\ell\)-haces constructibles}

\paragraph{1.4.1. \(A\)-haces constructibles (\(A\), \(\mathbb Z_\ell\)-álgebra finita).}
Sea \(A\) una \(\mathbb Z_\ell\)-álgebra finita. Se define una categoría, denotada por
\[
A\text{-}\fc(X),
\]
cuyos objetos se llaman \(A\)-haces constructibles, de la manera siguiente.

\begin{enumerate}[label=(\roman*)]
\item Un objeto de \(A\text{-}\fc(X)\) es un par \((F,u)\) formado por un \(\mathbb Z_\ell\)-haz constructible \(F\) y un morfismo de \(\mathbb Z_\ell\)-álgebras:
\[
u:A\longrightarrow \End(F).
\]
\item Si \((F,u)\) y \((F',u')\) son dos \(A\)-haces constructibles, una flecha \((F,u)\to(F',u')\) es un morfismo de \(\mathbb Z_\ell\)-haces
\[
f:F\longrightarrow F'
\]
tal que, para todo elemento \(a\) de \(A\), el diagrama que aparece a continuación sea conmutativo:
\[
\begin{tikzcd}
F \arrow[r,"f"] \arrow[d,"u(a)"'] & F' \arrow[d,"u'(a)"]\\
F \arrow[r,"f"'] & F'.
\end{tikzcd}
\]
\end{enumerate}

Se llamará \(A\)-haz constante torcido constructible a un \(A\)-haz constructible cuyo \(\mathbb Z_\ell\)-haz constructible subyacente sea constante torcido, y categoría de los \(A\)-haces constantes torcidos constructibles a la subcategoría plena de la categoría de los \(A\)-haces constructibles generada por los \(A\)-haces constantes torcidos constructibles.

Se sigue inmediatamente de (1.2.4) que, si \(X\) es conexo, la categoría de los \(A\)-haces constantes torcidos constructibles es equivalente a la categoría de los \(A\)-módulos de tipo finito dotados de una acción continua y \(A\)-lineal de \(\pi_1(X)\) sobre el \(\mathbb Z_\ell\)-módulo de tipo finito subyacente (1.2.4).
